\chapter{Gamma scalping et volatilité d'événement}\label{ch:s2:gamma-scalping-and-event-volatility}

Avant les annonces de résultats, les options sur une action intègrent un mouvement, puis l'action bouge plus ou moins que ce mouvement. Gao, Xing et Zhang ont établi que les straddles à la monnaie
sur des actions individuelles perdent en général de l'argent, mais que les mêmes straddles achetés trois jours avant l'annonce de résultats et conservés jusqu'à la
date de l'annonce ont rapporté en moyenne 3.34~\% de façon très significative. Détenir un straddle et le couvrir rapporte la différence entre le \omterm{def:s2:gamma-scalping-and-event-volatility:move}{mouvement implicite} et le
mouvement réalisé, diminuée du coût de la couverture. Ce chapitre introduit des sauts de résultats dans des actions synthétiques dont la variance d'événement implicite est trop faible
d'un montant contrôlé. Sur ces actions, le \omterm{def:s2:gamma-scalping-and-event-volatility:straddle}{straddle d'événement} rapporte 2.1~\% de sa prime par événement, soit un ratio de Sharpe d'environ 1 pour un book de 200 événements par an.
Un seul événement ne dit pratiquement rien, cependant~: l'écart type du rendement représente 29~\% de la prime. La fréquence de couverture qui donne les meilleurs résultats après coûts
est une fois par jour. La construction s'appelle \texttt{firm.eventvol}.

\section{Gamma scalping comme pari sur la volatilité réalisée contre la volatilité implicite}

Une option longue couverte en delta gagne, sur un intervalle court, environ $\tfrac12\Gamma S^2(\sigma_r^2 - \sigma_i^2)\,\mathrm{d}t$~: le gamma multiplié par
l'écart entre la variance réalisée et la variance à laquelle l'option était évaluée (livre~5, chapitre~25). Chaque couverture vend l'action après qu'elle a monté
et l'achète après qu'elle a baissé. Les transactions encaissent la variance du chemin, et le theta rémunère la variance implicite. Le trade est un pari sur la volatilité réalisée
contre la volatilité implicite, et la couverture en delta est le moyen de l'encaisser.

Les actions synthétiques de ce chapitre ont une volatilité diffusive voisine de 30~\%, et leur variance implicite comporte une prime de 25~\% par rapport à celle-ci. Acheter un
straddle à la monnaie à 21 jours et le couvrir pendant 20 jours de négociation, sans événement, fait perdre la prime, quelle que soit la couverture~:

\begin{center}
\small
\setlength{\tabcolsep}{4pt}%
\begin{tabular}{@{}p{4.6cm}>{\centering\arraybackslash}p{3.6cm}>{\centering\arraybackslash}p{2.4cm}>{\centering\arraybackslash}p{2.8cm}@{}}
\toprule
2\,000 straddles, sans événement & rendement moyen (\% de la prime) & écart type (\%) & coût de couverture (\%)\\
\midrule
jamais couvert & $-12.1$ & 62.4 & 0.0\\
chaque jour & $-12.4$ & 15.3 & 0.7\\
deux fois par jour & $-13.0$ & 11.5 & 1.0\\
toutes les deux heures & $-13.4$ & 8.5 & 1.5\\
toutes les heures & $-14.1$ & 6.7 & 2.1\\
quand le delta bouge de 0.05 & $-13.7$ & 7.2 & 1.6\\
quand le delta bouge de 0.10 & $-13.0$ & 8.6 & 1.2\\
quand le delta bouge de 0.20 & $-12.6$ & 12.6 & 0.7\\
\bottomrule
\end{tabular}
\end{center}

La couverture ne modifie pas ce que le trade espère rapporter. Ce rendement espéré est fixé par la variance implicite comparée à la variance réalisée. Ce que la couverture change, c'est la régularité avec laquelle le
trade le rapporte. Non couvert, le straddle est un pari directionnel dont l'écart type atteint 62~\% de sa prime. Couvert toutes les heures, l'écart type tombe à
6.7~\%, et 2.1~\% de la prime part dans les transactions sur l'action qui la couvrent. Les couvertures coûtent deux points de base de l'action négociée, une hypothèse de ce
chapitre.

Une règle de bande ne couvre que lorsque le delta a varié de plus d'un seuil depuis la dernière couverture, de sorte qu'elle négocie quand le chemin bouge et non quand
l'horloge sonne (\cref{fig:s2:gamma-scalping-and-event-volatility:scalping}). À coûts comparables, elle obtient un écart type plus faible que l'horloge~: une bande de
0.10 coûte 1.2~\% pour 8.6~\%, là où couvrir toutes les deux heures coûte 1.5~\% pour 8.5~\%.

\begin{omfigure}
\begin{tikzpicture}\pgfkeys{/omfr/.cd,daily/.initial={chaque jour},twice a day/.initial={deux fois par jour},every two hours/.initial={toutes les deux heures},hourly/.initial={toutes les heures},band 0.05/.initial={bande 0.05},band 0.10/.initial={bande 0.10},band 0.20/.initial={bande 0.20}}
\begin{axis}[width=10cm,height=5.6cm,xlabel={\small coût de couverture (\% de la prime)},ylabel={\small écart type du rendement (\% de la prime)},
  tick label style={font=\footnotesize},xmin=0.2,xmax=2.9,ymin=5,ymax=17,grid=major,grid style={gray!25},table/col sep=comma,
  nodes near coords={\pgfkeysvalueof{/omfr/\pgfplotspointmeta}},point meta=explicit symbolic,every node near coord/.append style={font=\scriptsize}]
\addplot[only marks,mark=*,omThm,every node near coord/.append style={anchor=west}] table[x=cost,y=sd,meta=rule]{figdata/strategies-2/04-gamma-scalping-and-event-volatility/scalping_clock.csv};
\addplot[only marks,mark=square*,omDef,every node near coord/.append style={anchor=east}] table[x=cost,y=sd,meta=rule]{figdata/strategies-2/04-gamma-scalping-and-event-volatility/scalping_band.csv};
\end{axis}
\end{tikzpicture}

\omcaption{Couverture d'un straddle long à 21 jours pendant 20 jours sur les actions synthétiques, sans événement~: coût de couverture de chaque règle en fonction de l'écart type
du rendement, exprimés l'un et l'autre en part de la prime. Règles d'horloge (ronds, de chaque jour à chaque heure) et
règles de bande (carrés, variations de delta de 0.05 à 0.20). Données~:
\texttt{s2\_eventvol.hedging(False)}.}\label{fig:s2:gamma-scalping-and-event-volatility:scalping}
\end{omfigure}

\section{Mouvements implicites autour des événements}

Une annonce de résultats tombe un jour connu et ajoute une masse de variance ce jour-là. Une échéance postérieure à l'événement la porte~; une échéance antérieure ne la porte pas.
Pour une échéance $\tau$ postérieure à l'événement, avec une volatilité implicite diffusive $\sigma_d$ et une variance d'événement $E$,
\[
\sigma^2(\tau)\,\tau = \sigma_d^2\,\tau + E,
\]
ainsi la variance d'événement relève davantage les volatilités des échéances courtes que celles des échéances longues, et deux échéances postérieures à l'événement suffisent pour résoudre à la fois $\sigma_d$ et
$E$. Dubinsky, Johannes, Kaeck et Seeger ont construit sur cette idée des modèles à forme réduite et des estimateurs afin de séparer l'incertitude portant sur les annonces de résultats
de la volatilité au jour le jour. L'incertitude anticipée était quantitativement importante, variait dans le temps et apportait de l'information sur la
volatilité future.

\begin{definition}[Mouvement implicite]\label{def:s2:gamma-scalping-and-event-volatility:move}
Le \emph{mouvement implicite}\index{mouvement implicite} d'un événement programmé est le rendement absolu espéré le jour de l'événement que les prix des options impliquent. Sous un
rendement normal de l'événement, avec la variance d'événement $E$ extraite de la structure par terme, il vaut $\sqrt{2/\pi}\,\sqrt{E}$. Les traders l'approchent souvent par
le prix du straddle à la monnaie arrivant à échéance juste après l'événement, en proportion du prix de l'action.
\end{definition}

Sur l'une des actions du chapitre, avec une variance implicite diffusive et une variance d'événement implicite de 0.00135, une échéance à 5 jours se négocie à 39.5~\% et une
échéance à 21 jours à 32.3~\%. Les deux volatilités restituent exactement la variance d'événement (\texttt{event\_variance})~: un écart type d'événement de 3.67~\%.

\begin{definition}[Straddle d'événement]\label{def:s2:gamma-scalping-and-event-volatility:straddle}
Un \emph{straddle d'événement}\index{straddle d'événement} est un straddle à la monnaie acheté peu avant un événement programmé et vendu peu après, en général
couvert en delta, qui rapporte lorsque le mouvement réalisé lors de l'événement dépasse le \omterm{def:s2:gamma-scalping-and-event-volatility:move}{mouvement implicite} d'un montant supérieur à la décote temporelle et aux coûts de détention.
\end{definition}

Le marché synthétique introduit une mauvaise évaluation. Le mouvement réel de chaque événement a un écart type voisin de 5~\%, tiré événement par événement. Le marché affiche une
variance d'événement implicite égale à la vraie multipliée par 0.70, puis par une erreur lognormale~: le \omterm{def:s2:gamma-scalping-and-event-volatility:move}{mouvement implicite} est en moyenne trop faible, et erroné d'un montant différent
à chaque fois. Sur les 2\,000 événements, l'écart absolu moyen était de 4.35~\% et le \omterm{def:s2:gamma-scalping-and-event-volatility:move}{mouvement implicite} moyen de 3.62~\%. Le \omterm{def:s2:gamma-scalping-and-event-volatility:move}{mouvement implicite} médian valait 0.81 du
mouvement réel. Le book achète le straddle à 21 jours à la clôture trois jours avant l'événement, le couvre selon une règle et le vend à la clôture du jour de l'événement.
Les options coûtent 2~\% de la prime pour l'aller-retour, et les couvertures deux points de base.

\begin{center}
\small
\begin{tabular}{@{}p{3.9cm}>{\centering\arraybackslash}p{2.0cm}>{\centering\arraybackslash}p{2.0cm}>{\centering\arraybackslash}p{2.0cm}>{\centering\arraybackslash}p{2.0cm}>{\centering\arraybackslash}p{2.0cm}@{}}
\toprule
2\,000 événements & moyenne (\%) & écart type (\%) & moyenne\slash écart type & coût de couverture (\%) & 200 événements par an\\
\midrule
jamais couvert & 2.0 & 31 & 0.064 & 0.00 & 0.91\\
chaque jour & 2.1 & 29 & 0.073 & 0.06 & 1.04\\
deux fois par jour & 2.0 & 28 & 0.071 & 0.18 & 1.00\\
toutes les deux heures & 1.8 & 27 & 0.064 & 0.23 & 0.91\\
toutes les heures & 1.7 & 27 & 0.063 & 0.30 & 0.89\\
quand le delta bouge de 0.05 & 1.7 & 27 & 0.063 & 0.24 & 0.89\\
quand le delta bouge de 0.10 & 2.0 & 28 & 0.070 & 0.18 & 1.00\\
quand le delta bouge de 0.20 & 1.9 & 28 & 0.068 & 0.13 & 0.96\\
\bottomrule
\end{tabular}
\end{center}

La dernière colonne multiplie le ratio par événement par $\sqrt{200}$, soit un portefeuille de cinquante noms avec quatre publications chacun, et suppose que les événements sont indépendants,
ce qui est le cas dans cette simulation et ne l'est pas sur un marché. Le saut de cours intervient à l'ouverture, avant que quiconque puisse couvrir, si bien qu'aucune règle de couverture ne le capture. La couverture
ne fait que réduire le bruit des journées qui l'entourent. La couverture quotidienne donne les meilleurs résultats une fois les coûts déduits~: couvrir plus souvent réduit l'écart-type de 29~\% à 27~\% et
coûte plus qu'elle ne rapporte.

Reproduire l'effet compte davantage que n'importe quelle règle. Avec une variance d'événement implicite juste en moyenne, le portefeuille couvert chaque jour perd 2.2~\% par événement. Lorsqu'elle est
sous-évaluée de 15~\%, il perd 0.1~\%. Ce n'est qu'à 30~\% de sous-évaluation qu'il gagne ses 2.1~\%. Le straddle doit d'abord absorber son propre theta, la prime diffusive et ses coûts
avant que l'événement ne rapporte quoi que ce soit.

\section{L'effondrement de la volatilité après les événements}

\begin{definition}[Effondrement de la volatilité]\label{def:s2:gamma-scalping-and-event-volatility:crush}
L'\emph{effondrement de la volatilité}\index{effondrement de la volatilité} (volatility crush) est la baisse de la volatilité implicite après un événement programmé, lorsque la variance de l'événement sort
de la durée de vie restante des options~; elle est maximale pour les échéances les plus courtes, celles dont l'événement avait le plus relevé la volatilité.
\end{definition}

Sur les actions synthétiques, la volatilité implicite médiane à 21 jours valait 36.6~\% à l'achat du straddle et 32.0~\% après l'événement, soit une baisse médiane de 8.6~\%. Le
détenteur de straddle perd ce vega le jour de l'événement quoi que fasse l'action, si bien que le saut de cours doit le compenser.

Le même schéma apparaît sur l'indice autour des communiqués programmés de la Réserve fédérale. Sur les 125 jours de communiqué du FOMC programmés entre janvier 2011 et
septembre 2026, le rendement logarithmique absolu moyen du S\&P 500 valait 0.83~\%, contre 0.71~\% les autres jours. Son rendement quadratique moyen était 1.146 fois plus élevé,
soit un mouvement propre à l'événement d'écart-type d'environ 0.41~\%. Le VIX à 9 jours de Cboe a baissé lors de 65.6~\% des jours de communiqué, contre 53.1~\% des autres jours, avec une baisse moyenne de
3.05~\% (variation logarithmique), contre une hausse de 0.09~\% les autres jours. Le VIX à 30 jours, que le communiqué relève moins, a baissé de 1.10~\% en moyenne.
Un événement macroéconomique est modeste à côté d'une publication de résultats, et l'effondrement qu'il laisse dans la volatilité la plus courte reste visible.

\section{Fréquence et chemin de la couverture}

L'erreur sur le \omterm{def:s2:gamma-scalping-and-event-volatility:move}{mouvement implicite} détermine chaque événement (la \cref{fig:s2:gamma-scalping-and-event-volatility:error}). Classés en quintiles selon l'erreur, soit le logarithme du \omterm{def:s2:gamma-scalping-and-event-volatility:move}{mouvement implicite} rapporté au mouvement réel, le
straddle couvert quotidiennement a rapporté 8.0~\% de sa prime lorsque le \omterm{def:s2:gamma-scalping-and-event-volatility:move}{mouvement implicite} était sous-évalué de 0.46 en logarithme. Là où le \omterm{def:s2:gamma-scalping-and-event-volatility:move}{mouvement implicite}
était à peu près juste (erreur de 0.05), il a perdu 4.5~\%. La corrélation entre l'erreur et le rendement, événement par événement, n'est que de $-0.13$, car le
mouvement observé est un tirage unique~: un événement bien évalué peut encore s'ouvrir sur un saut de trois écarts-types.

\begin{omfigure}
\begin{tikzpicture}
\begin{axis}[width=10cm,height=5.4cm,ybar,bar width=16pt,xlabel={\small quintile d'événements selon l'erreur du mouvement implicite (1~: le plus sous-évalué)},
  ylabel={\small rendement moyen (\% de la prime)},tick label style={font=\footnotesize},xmin=0.4,xmax=5.6,ymin=-6,ymax=10,xtick={1,2,3,4,5},
  grid=major,grid style={gray!25},table/col sep=comma]
\addplot[fill=omThm!70,draw=omThm] table[x=quintile,y=mean]{figdata/strategies-2/04-gamma-scalping-and-event-volatility/error.csv};
\end{axis}
\end{tikzpicture}

\omcaption{Rendement moyen du \omterm{def:s2:gamma-scalping-and-event-volatility:straddle}{straddle d'événement} couvert quotidiennement, par quintile des 2\,000 événements synthétiques, classés selon l'erreur du \omterm{def:s2:gamma-scalping-and-event-volatility:move}{mouvement implicite} (logarithme du \omterm{def:s2:gamma-scalping-and-event-volatility:move}{mouvement implicite}
sur le mouvement réel~; erreurs moyennes $-0.46$, $-0.30$, $-0.21$, $-0.11$ et 0.05). Données~: \texttt{s2\_eventvol.by\_error}.}\label{fig:s2:gamma-scalping-and-event-volatility:error}
\end{omfigure}

Le chemin compte autant que la fréquence. Une règle calendaire couvre les journées calmes comme les journées agitées, tandis qu'une règle par bandes attend que le delta bouge. Autour
d'un événement, le chemin a une structure~: les journées diffusives qui le précèdent, puis le saut de cours à l'ouverture. Ce saut ne peut pas être couvert, si bien que toute règle qui couvre les journées
qui précèdent le saut paie pour réduire une variance qu'elle aurait pu laisser de côté. C'est pourquoi la règle quotidienne et la bande à 0.10 font jeu égal, et pourquoi la couverture horaire perd 0.4~\%
de la prime par rapport à la règle quotidienne sur trois jours.

\section{Fiches de stratégie}

\begin{strategyfile}{Straddle d'événement sur résultats}\label{strat:s2:gamma-scalping-and-event-volatility:earnings}
\sfield{Qui paie la stratégie, et pourquoi} Vendeurs d'options qui sous-estiment l'incertitude d'une annonce, d'autant plus lorsque les signaux de l'entreprise sont bruités et que ses
options sont coûteuses à échanger.
\sfield{Instruments et plateformes} Options cotées sur action unique, straddles à la monnaie avec la première échéance après la publication.
\sfield{Signal} \omterm{def:s2:gamma-scalping-and-event-volatility:move}{Mouvement implicite} comparé à une prévision du mouvement~: surprises passées, kurtosis, taille, dispersion des analystes.
\sfield{Dimensionnement et exécution} De nombreuses petites positions sur plusieurs noms et saisons de publication~; achetées quelques jours à l'avance, vendues après la publication.
\sfield{Coûts} Fourchettes des options, les plus larges là où l'effet est le plus fort.
\sfield{Comment elle s'éteint} Encombrement (crowding) qui fait monter les \omterm{def:s2:gamma-scalping-and-event-volatility:move}{mouvements implicites}~; coûts qui absorbent la prime.
\sfield{Horizon, capacité, infrastructure} Quelques jours par transaction~; un calendrier des résultats et des surfaces d'options pour des centaines de noms.
\sfield{Backtester honnêtement} Fourchettes affichées, et non prix milieu~; dates d'annonce telles qu'elles étaient connues alors~; l'effondrement dans le prix de sortie.
\sfield{Sources} Gao, Xing et Zhang (2018)~: 3.34~\% en moyenne de trois jours avant à la date d'annonce~; Dubinsky, Johannes, Kaeck et Seeger
(2019)~; ce chapitre~: 2.1~\% par événement, soit environ 1.0 pour 200 événements par an.
\end{strategyfile}

\begin{strategyfile}{Vente de volatilité après l'événement}\label{strat:s2:gamma-scalping-and-event-volatility:postevent}
\sfield{Qui paie la stratégie, et pourquoi} Détenteurs d'options d'événement qui vendent après l'événement, et la prime de variance diffusive qui subsiste.
\sfield{Instruments et plateformes} Options sur action unique ou sur indice après la publication des résultats ou le communiqué.
\sfield{Signal} Volatilité implicite restante comparée à la volatilité diffusive prévue une fois l'événement passé.
\sfield{Dimensionnement et exécution} Straddles vendus couverts en delta ou variance, dimensionnés par un scénario de stress de saut de cours.
\sfield{Coûts} Fourchettes des options~; couverture.
\sfield{Comment elle s'éteint} Nouvelles qui en découlent (prévisions, retraitements) et apportent un second événement.
\sfield{Horizon, capacité, infrastructure} De quelques jours à quelques semaines~; surfaces d'options mises à jour après chaque événement.
\sfield{Backtester honnêtement} Sortie et entrée aux cours affichés après l'événement, et non à la clôture précédente.
\sfield{Sources} Aucun chiffre de performance vérifié.
\end{strategyfile}

\begin{strategyfile}{Gamma scalping avec bandes de couverture}\label{strat:s2:gamma-scalping-and-event-volatility:bands}
\sfield{Qui paie la stratégie, et pourquoi} Vendeurs d'options qui cotent la volatilité en dessous de ce que livre le chemin.
\sfield{Instruments et plateformes} Options achetées sur sous-jacents liquides~; l'action ou le contrat à terme pour la couverture.
\sfield{Signal} Volatilité réalisée prévue comparée à la volatilité implicite.
\sfield{Dimensionnement et exécution} Bandes de delta fixées en arbitrant coûts et variance~; couverture sur les mouvements, et non à l'horloge.
\sfield{Coûts} Négociation de l'action ou des contrats à terme à chaque couverture.
\sfield{Comment elle s'éteint} Volatilité implicite au-dessus de la réalisée en moyenne (la prime de variance)~: un portefeuille long en gamma la perd.
\sfield{Horizon, capacité, infrastructure} De quelques jours à quelques semaines~; un moteur de couverture avec des deltas en temps réel.
\sfield{Backtester honnêtement} Trajectoires intrajournalières~; coûts de couverture par transaction.
\sfield{Sources} Ce chapitre~: sans événement, le straddle long perd 12~\% à 14~\% de sa prime quelle que soit la règle~; la couverture fixe le niveau de bruit.
\end{strategyfile}

\begin{strategyfile}{Variance d'événements macroéconomiques}\label{strat:s2:gamma-scalping-and-event-volatility:macro}
\sfield{Qui paie la stratégie, et pourquoi} Acheteurs de protection autour des décisions des banques centrales et des publications de données, et de l'effondrement qui suit.
\sfield{Instruments et plateformes} Options sur indice à courte échéance et échéances hebdomadaires autour des jours de communiqué.
\sfield{Signal} La variance d'événement dans les échéances courtes, comparée aux mouvements passés de l'indice ces jours-là.
\sfield{Dimensionnement et exécution} Taille réduite~; dimensionnée selon les plus forts mouvements passés les jours d'événement.
\sfield{Coûts} Fourchettes des options sur indice à courte échéance.
\sfield{Comment elle s'éteint} Une décision surprise~; une réunion non programmée.
\sfield{Horizon, capacité, infrastructure} Quelques jours~; un calendrier des événements programmés.
\sfield{Backtester honnêtement} Dates programmées telles qu'elles étaient publiées alors~; réunions non programmées exclues du signal et conservées dans le risque.
\sfield{Sources} Calendriers de la Federal Reserve et historiques de Cboe (ce chapitre~: 125 jours de communiqué, VIX9D en baisse de 3.05~\% en moyenne).
\end{strategyfile}

\section{Tutoriel~: le mouvement évalué et le mouvement observé}\label{tut:s2:gamma-scalping-and-event-volatility:tut}

\begin{tutorial}
\textbf{Objectif.} Simuler des fenêtres d'événement, évaluer et couvrir des \omterm{def:s2:gamma-scalping-and-event-volatility:straddle}{straddles d'événement} selon plusieurs règles, et mesurer l'\omterm{def:s2:gamma-scalping-and-event-volatility:crush}{effondrement de la volatilité} et les statistiques des jours FOMC.
\textbf{État final~:} les deux tableaux et les deux figures.
\begin{tutsteps}
\item \textbf{Le straddle et ses couvertures}.
\omcode{code/firm/eventvol/firm_eventvol.py}{87}{118}{Un straddle long par fenêtre, couvert en delta selon une règle calendaire ou par bandes, avec les coûts.}\label{lst:s2:gamma-scalping-and-event-volatility:pnl}
\item \textbf{Les règles de couverture}.
\omcode{code/strategies-2/04-gamma-scalping-and-event-volatility/python/s2_eventvol.py}{34}{51}{Rendement, bruit et coût de chaque règle.}\label{lst:s2:gamma-scalping-and-event-volatility:rules}
\item \textbf{Exécuter} \texttt{s2\_fetch\_fomc.py} une fois, puis \texttt{hedging()}, \texttt{hedging(False)}, \texttt{by\_error()}, \texttt{crush()} et
\texttt{fig\_eventvol.py}.
\end{tutsteps}
\textbf{Pour aller plus loin.} Faire dépendre l'erreur sur le \omterm{def:s2:gamma-scalping-and-event-volatility:move}{mouvement implicite} d'un élément observable (surprises passées) et ne prendre position que sur le quintile le plus susceptible
d'être sous-évalué~; rendre les événements corrélés entre les noms d'une saison~; ajouter un saut de cours à l'ouverture les jours sans événement.
\end{tutorial}

\section{Construction~: volatilité d'événement}\label{bld:s2:gamma-scalping-and-event-volatility:eventvol}

\begin{build}
\bfield{Rôle} Variance d'événement issue de la structure par terme, \omterm{def:s2:gamma-scalping-and-event-volatility:move}{mouvements implicites}, fenêtres d'événement avec sauts de cours à l'ouverture, et straddles couverts en temps discret.
\bfield{Interface} \texttt{EventConfig(\dots)}, \texttt{event\_variance(v1, t1, v2, t2)}, \texttt{implied\_move(event\_var)},
\texttt{windows(cfg, days, event)}, \texttt{straddle\_pnl(w, cfg, every, band)}.
\bfield{Règles} Le saut a lieu à l'ouverture du jour de l'événement, entre deux couvertures~; la variance d'événement quitte la volatilité implicite après le saut~; des coûts s'appliquent à chaque
couverture et aux options.
\bfield{Tests d'acceptation} \texttt{code/firm/eventvol/tests/}~: la variance d'événement retrouvée exactement à partir de deux échéances~; l'\omterm{def:s2:gamma-scalping-and-event-volatility:crush}{effondrement de la volatilité} après chaque événement et
aucun sans événement~; sur un chemin immobile, seuls la décroissance temporelle et les coûts.
\bfield{Pour aller plus loin} Plusieurs événements par fenêtre~; sauts les jours sans événement~; volatilité diffusive stochastique.
\end{build}

\begin{omsources}
\item A.~Dubinsky, M.~Johannes, A.~Kaeck et N.~J.~Seeger, ``Option pricing of earnings announcement risks'', \emph{Review of Financial Studies}
32(2), 2019.
\item C.~Gao, Y.~Xing et X.~Zhang, ``Anticipating uncertainty: straddles around earnings announcements'', \emph{Journal of Financial and Quantitative
Analysis} 53(6), 2018.
\item Federal Reserve, calendriers des réunions du FOMC, 2011--2026.
\item Cboe Global Markets, historiques quotidiens de SPX, VIX et VIX9D.
\end{omsources}

\section{Exercices}

\begin{exercise}[$\star$]\label{exo:s2:gamma-scalping-and-event-volatility:1}
La variance d'un événement vaut 0.0025. Calculer son \omterm{def:s2:gamma-scalping-and-event-volatility:move}{mouvement implicite}.
\end{exercise}

\begin{exercise}[$\star$]\label{exo:s2:gamma-scalping-and-event-volatility:2}
La volatilité implicite diffusive d'une action vaut 30~\% et un événement ajoute une variance de 0.0025. Calculer la volatilité implicite d'une échéance située 10 jours de bourse plus tard, après l'événement.
\end{exercise}

\begin{exercise}[$\star$]\label{exo:s2:gamma-scalping-and-event-volatility:3}
Expliquer pourquoi aucune règle de couverture ne peut capter un saut de cours lié à des résultats.
\end{exercise}

\begin{exercise}[$\star\star$]\label{exo:s2:gamma-scalping-and-event-volatility:4}
Gao, Xing et Zhang ont montré que les straddles perdent en général, mais que les \omterm{def:s2:gamma-scalping-and-event-volatility:straddle}{straddles d'événement} gagnent. Concilier les deux résultats.
\end{exercise}

\begin{exercise}[$\star\star$]\label{exo:s2:gamma-scalping-and-event-volatility:5}
Expliquer pourquoi une règle à bandes subit moins de bruit, à coût égal, qu'une règle calendaire.
\end{exercise}

\begin{exercise}[$\star\star$]\label{exo:s2:gamma-scalping-and-event-volatility:6}
La moyenne des carrés des rendements du S\&P 500 les jours de réunion du FOMC vaut 1.146 fois celle des autres jours. Pourquoi le pari sur l'événement macroéconomique est-il modeste à côté du pari sur les résultats~?
\end{exercise}

\begin{exercise}[$\star\star\star$]\label{exo:s2:gamma-scalping-and-event-volatility:7}
\emph{Programmation.} Exécuter \texttt{hedging(True, 0.0)} et \texttt{hedging(True, -0.15)}. De combien l'événement doit-il être mal évalué pour que le straddle
couvert quotidiennement atteigne le seuil de rentabilité, et pourquoi la réponse n'est-elle pas nulle~?
\end{exercise}

\begin{exercise}[$\star\star\star$]\label{exo:s2:gamma-scalping-and-event-volatility:8}
\emph{Trouver la faille.} «~Nos \omterm{def:s2:gamma-scalping-and-event-volatility:straddle}{straddles d'événement} rapportent 2.1~\% par événement avec un ratio de Sharpe de 1.04 par an, donc nous pouvons dimensionner le livre sur ce ratio de Sharpe.~»
\end{exercise}

\section{Problème~: le mouvement évalué et le mouvement réalisé}

\begin{problem}[{Problème du week-end --- négocier des événements}]\label{pb:s2:gamma-scalping-and-event-volatility:1}
Les événements synthétiques du chapitre, l'historique des réunions du FOMC et les travaux de recherche publics.

\medskip\noindent\textbf{Partie I --- Gamma.}
\begin{enumerate}
\item Écrire le P\&L de l'option couverte sur un court intervalle de temps et l'expliquer.
\item Pourquoi la couverture ne modifie-t-elle pas le gain espéré d'un straddle long~?
\item Combien a rapporté le straddle sans événement selon chaque règle de couverture~?
\item Définir une règle à bandes et la comparer à une règle calendaire.
\end{enumerate}

\medskip\noindent\textbf{Partie II --- Événements.}
\begin{enumerate}[resume]
\item Comment extrait-on la variance d'événement de la structure par terme~?
\item Définir le \omterm{def:s2:gamma-scalping-and-event-volatility:move}{mouvement implicite} et le \omterm{def:s2:gamma-scalping-and-event-volatility:straddle}{straddle d'événement}.
\item Qu'ont trouvé Dubinsky et ses coauteurs, ainsi que Gao, Xing et Zhang~?
\item Comment le marché synthétique évalue-t-il mal les événements~?
\end{enumerate}

\medskip\noindent\textbf{Partie III --- L'\omterm{def:s2:gamma-scalping-and-event-volatility:crush}{effondrement de la volatilité}.}
\begin{enumerate}[resume]
\item Définir l'\omterm{def:s2:gamma-scalping-and-event-volatility:crush}{effondrement de la volatilité} et donner son ampleur sur les actions synthétiques.
\item Que se passe-t-il pour le S\&P 500 et pour le VIX9D les jours de réunion du FOMC~?
\item Pourquoi le VIX à 30 jours baisse-t-il moins que le VIX à 9 jours~?
\item Combien l'effondrement coûte-t-il au détenteur du straddle~?
\end{enumerate}

\medskip\noindent\textbf{Partie IV --- Le verdict.}
\begin{enumerate}[resume]
\item Énoncer le \emph{résultat nommé}~: le rendement du \omterm{def:s2:gamma-scalping-and-event-volatility:straddle}{straddle d'événement} par événement face à l'erreur sur le \omterm{def:s2:gamma-scalping-and-event-volatility:move}{mouvement implicite}, et la fréquence de couverture qui
la maximise après coûts.
\item Donner le rendement par quintile de l'erreur sur le \omterm{def:s2:gamma-scalping-and-event-volatility:move}{mouvement implicite}.
\item Pourquoi la corrélation événement par événement n'est-elle que de $-0.13$~?
\item De combien les événements doivent-ils être mal évalués pour que le pari rapporte~?
\item Que suppose le ratio de Sharpe calculé sur 200 événements par an~?
\item Comment backtester honnêtement un livre de straddles sur résultats~?
\item Quelle fiche de stratégie est le plus proche de la prime de variance du chapitre~1~?
\item En une phrase~: qu'achète un \omterm{def:s2:gamma-scalping-and-event-volatility:straddle}{straddle d'événement}~?
\end{enumerate}
\end{problem}

\section{Questions d'entretien}

\begin{interviewq}[$\star$ \iqroles{trader}]\label{iq:s2:gamma-scalping-and-event-volatility:1}
Qu'est-ce que le gamma scalping, et quand rapporte-t-il~?
\end{interviewq}

\begin{interviewq}[$\star\star$ \iqroles{trader}]\label{iq:s2:gamma-scalping-and-event-volatility:2}
Une action publie ses résultats demain. Comment lisez-vous le \omterm{def:s2:gamma-scalping-and-event-volatility:move}{mouvement implicite} à partir de la chaîne d'options~?
\end{interviewq}

\begin{interviewq}[$\star\star$ \iqroles{chercheur}]\label{iq:s2:gamma-scalping-and-event-volatility:3}
Comment testeriez-vous si les mouvements liés aux résultats sont sous-évalués~?
\end{interviewq}

\begin{interviewq}[$\star\star$ \iqroles{risque}]\label{iq:s2:gamma-scalping-and-event-volatility:4}
Un livre est vendeur de \omterm{def:s2:gamma-scalping-and-event-volatility:straddle}{straddles d'événement} sur 200 noms pendant la saison des résultats. Quel est son risque, et comment le soumettriez-vous à un stress~?
\end{interviewq}

\begin{interviewq}[$\star\star$ \iqroles{développeur}]\label{iq:s2:gamma-scalping-and-event-volatility:5}
Concevez un moteur de couverture en delta avec des règles à bandes. Que doit-il connaître, et que peut-il mal tourner à l'ouverture~?
\end{interviewq}

\begin{interviewq}[$\star\star\star$ \iqroles{chercheur}]\label{iq:s2:gamma-scalping-and-event-volatility:6}
Deux échéances après un événement, $\tau_1 < \tau_2$, ont des volatilités implicites $\sigma_1$ et $\sigma_2$. Établir la variance d'événement et la volatilité diffusive, et
dire ce qui ne va pas si une échéance tombe avant l'événement.
\end{interviewq}
